% Chunk: Chapter 14, Section 14.1, The Dirac Equation.
% Source: Weinberg QFT Vol. I, physical pp. 588--595 (printed pp. 565--572).
% Coverage: equations (14.1.1)--(14.1.47); one body footnote; citations [1]--[3].
% Boundary: begins at the Section 14.1 heading on physical p. 588 and ends at
% Eq. (14.1.47) on physical p. 595, immediately before the Section 14.2 heading.
% Source anomaly retained: both Fourier-transform displays on pp. 571--572
% visibly print d^3p although their integrands are functions of x.
% Uncertainty: none.

\subsection{The Dirac Equation}

We shall limit ourselves in this chapter to problems in which bound states
arise because of the Coulomb interaction of electrons (or muons) with heavy
charged particles such as atomic nuclei. As shown in Section 13.6, this
interaction may be taken into account by adding to the interaction Lagrangian
a term\footnote{In this chapter we return to the use of an upper case $\Psi$ to
denote the electron field in the Heisenberg picture, reserving a lower case
$\psi$ for the Dirac field with time-dependence governed solely by the c-number
external field $\mathcal A^\mu(x)$.} representing the effects of a c-number
external vector potential $\mathcal A^\mu(x)$:
\begin{equation}
\begin{aligned}
\mathcal L_{\mathcal A}={}&-ie\bar\Psi\gamma^\mu\Psi\mathcal A_\mu
-\frac14(Z_3-1)(\partial^\mu\mathcal A^\nu-\partial^\nu\mathcal A^\mu)
  (\partial_\mu A_\nu-\partial_\nu A_\mu)\\
&-ie(Z_2-1)\mathcal A_\mu\bar\Psi\gamma^\mu\Psi,
\end{aligned}
\tag{14.1.1}\label{eq:14.1.1}
\end{equation}
which is obtained by replacing the quantum vector potential $A^\mu$ with
$A^\mu+\mathcal A^\mu$ in the interaction part of Eq.~\eqref{eq:11.1.6}. For
instance, for a single heavy particle of charge $Ze$ at the origin,
\begin{equation}
  \mathcal A^0(x)=\frac{Ze}{4\pi|\mathbf x|},\qquad \mathbf A(x)=0.
  \tag{14.1.2}\label{eq:14.1.2}
\end{equation}
It is the interaction \eqref{eq:14.1.1} that must be taken into account to all
orders. In this section we will consider the theory with only this interaction,
leaving radiative corrections to subsequent sections.

Physicists learn in kindergarten to approach this sort of problem by solving
the wave equation of Dirac in the presence of the external field. It might seem
unnecessary to derive this equation here, but as emphasized in Chapter 1,
Dirac's original motivation for this equation as a sort of relativistic
Schr\"odinger equation does not stand up to inspection. Also, in the course of
our derivation we will discover the normalization conditions that have to be
imposed on the solutions of the Dirac equation, which seemed somewhat
\emph{ad hoc} in Dirac's approach. The solutions of the Dirac equation
discussed here will be important ingredients in our treatment of radiative
corrections in the next section.

We will work here in a version of the Heisenberg picture, in which the
time-dependence of operators is determined by a Hamiltonian including the
external field interaction \eqref{eq:14.1.1}, but no other interactions. The
electron field $\psi(x)$ in this picture satisfies the field equation
\begin{equation}
  \left[\gamma^\lambda\frac{\partial}{\partial x^\lambda}
  +m+ie\gamma^\lambda\mathcal A_\lambda(x)\right]\psi(x)=0.
  \tag{14.1.3}\label{eq:14.1.3}
\end{equation}
This is not the Dirac equation in the original sense of
Dirac~\hyperlink{ref:14.1}{[1]}, because $\psi(x)$ here is not a c-number wave
function but a quantum operator. The c-number Dirac wave functions are defined
by
\begin{equation}
  u_N(x)\equiv\bra{\Omega}\psi(x)\ket{N},
  \tag{14.1.4}\label{eq:14.1.4}
\end{equation}
\begin{equation}
  v_N(x)\equiv\bra{N}\psi(x)\ket{\Omega},
  \tag{14.1.5}\label{eq:14.1.5}
\end{equation}
where the $\ket{N}$ are a complete orthonormal set of state-vectors, with
$\ket{\Omega}$ the vacuum. It follows immediately from Eq.~\eqref{eq:14.1.3}
that these functions satisfy the homogeneous Dirac equation
\begin{equation}
\begin{aligned}
\left[\gamma^\lambda\frac{\partial}{\partial x^\lambda}
 +m+ie\gamma^\lambda\mathcal A_\lambda(x)\right]u_N(x)
&=\left[\gamma^\lambda\frac{\partial}{\partial x^\lambda}
 +m+ie\gamma^\lambda\mathcal A_\lambda(x)\right]v_N(x)\\
&=0.
\end{aligned}
\tag{14.1.6}\label{eq:14.1.6}
\end{equation}
We can also derive a normalization condition from the equal-time
anticommutation relations for the Dirac field. These are unaffected by the
interaction \eqref{eq:14.1.1}, and therefore take the same form as for the free
fields:
\begin{equation}
  \{\psi(\mathbf x,t),\bar\psi(\mathbf y,t)\}
  =i\gamma^0\delta^3(\mathbf x-\mathbf y).
  \tag{14.1.7}\label{eq:14.1.7}
\end{equation}
Taking the vacuum expectation value and inserting a sum over the states
$\ket{N}$, we find
\begin{equation}
  \sum_Nu_N(\mathbf x,t)u_N^\dagger(\mathbf y,t)
  +\sum_Nv_N(\mathbf x,t)v_N^\dagger(\mathbf y,t)
  =\delta^3(\mathbf x-\mathbf y),
  \tag{14.1.8}\label{eq:14.1.8}
\end{equation}
it being understood that the sum over $N$ includes an integral over continuum
states as well as a sum over any discrete bound states.

We are chiefly interested in the case of a time-independent external field,
like \eqref{eq:14.1.2}. In such cases, the states $\ket{N}$ may be taken as
eigenstates of the Hamiltonian (including the interaction \eqref{eq:14.1.1})
with energies $E_N$. Time-translation invariance then tells us that the
$u_N(x)$ and $v_N(x)$ have the time-dependence:
\begin{equation}
  u_N(\mathbf x,t)=e^{-iE_Nt}u_N(\mathbf x),\qquad
  v_N(\mathbf x,t)=e^{+iE_Nt}v_N(\mathbf x).
  \tag{14.1.9}\label{eq:14.1.9}
\end{equation}
The homogeneous Dirac equations \eqref{eq:14.1.6} then become
\begin{equation}
  i\gamma^0\left[\boldsymbol\gamma\cdot\boldsymbol\nabla+m
  +ie\gamma^\lambda\mathcal A_\lambda(\mathbf x)\right]u_N(\mathbf x)
  =E_Nu_N(\mathbf x),
  \tag{14.1.10}\label{eq:14.1.10}
\end{equation}
\begin{equation}
  i\gamma^0\left[\boldsymbol\gamma\cdot\boldsymbol\nabla+m
  +ie\gamma^\lambda\mathcal A_\lambda(\mathbf x)\right]v_N(\mathbf x)
  =-E_Nv_N(\mathbf x).
  \tag{14.1.11}\label{eq:14.1.11}
\end{equation}
The minus sign on the right-hand side of Eq.~\eqref{eq:14.1.11} shows that the
$v_N$ are the famous `negative-energy' solutions of Dirac. As shown by
Eq.~\eqref{eq:14.1.8}, these negative-energy solutions are needed to make up a
complete set of wave functions. Of course, for moderate external fields there
are no negative-energy states in the theory, so all $E_N$ are positive, but
there is still an important difference between the states with non-vanishing
$u_N$ or $v_N$: the definitions \eqref{eq:14.1.4} and \eqref{eq:14.1.5} show
that a state can have $u_N\ne0$ or $v_N\ne0$ only if it has charge $-e$ or
$+e$, respectively. It is in this sense that negative-energy solutions of the
Dirac equation have something to do with the existence of antiparticles.
However, this argument has nothing to do with the details of the Dirac
equation, or even with the spin of the electron.

From the Dirac wave equations \eqref{eq:14.1.10} and \eqref{eq:14.1.11} we can
easily see that wave functions of different energy are orthogonal. That is,
\[
  (E_M-E_N^*)(u_N^\dagger u_M)
  =\boldsymbol\nabla\cdot(u_N^\dagger i\gamma^0\boldsymbol\gamma u_M),
\]
so if $|\mathbf x|^2(u_N^\dagger i\gamma^0\boldsymbol\gamma u_M)$ remains
bounded as $|\mathbf x|\to0$ and $|\mathbf x|\to\infty$, then
\begin{equation}
  \int d^3x\,(u_N^\dagger(\mathbf x)u_M(\mathbf x))=0
  \quad\text{if}\quad E_N\ne E_M^*.
  \tag{14.1.12}\label{eq:14.1.12}
\end{equation}
With similar boundary conditions for the $v_N$, we find in the same way that
\begin{equation}
  \int d^3x\,(v_N^\dagger(\mathbf x)v_M(\mathbf x))=0
  \quad\text{if}\quad E_N\ne E_M^*,
  \tag{14.1.13}\label{eq:14.1.13}
\end{equation}
\begin{equation}
  \int d^3x\,(u_N^\dagger(\mathbf x)v_M(\mathbf x))=0
  \quad\text{if}\quad E_N\ne-E_M^*.
  \tag{14.1.14}\label{eq:14.1.14}
\end{equation}
Taking $N=M$, Eqs.~\eqref{eq:14.1.12} and \eqref{eq:14.1.13} tell us that the
energies are all real. Dropping the complex conjugation of $E_M$ in
Eqs.~\eqref{eq:14.1.12}--\eqref{eq:14.1.14}, we see that $u$s of different
energy are orthogonal, $v$s of different energy are orthogonal, and (as long
as the potential is not strong enough to produce negative energy states) all
$u$s are orthogonal to all $v$s. By a suitable choice of the discrete quantum
numbers that characterize the states along with the energy, we can then always
arrange that
\begin{equation}
  \int d^3x\,(u_N^\dagger(\mathbf x)u_M(\mathbf x))=0\quad\text{if}\quad N\ne M,
  \tag{14.1.15}\label{eq:14.1.15}
\end{equation}
\begin{equation}
  \int d^3x\,(v_N^\dagger(\mathbf x)v_M(\mathbf x))=0\quad\text{if}\quad N\ne M,
  \tag{14.1.16}\label{eq:14.1.16}
\end{equation}
\begin{equation}
  \int d^3x\,(u_N^\dagger(\mathbf x)v_M(\mathbf x))=0.
  \tag{14.1.17}\label{eq:14.1.17}
\end{equation}

Multiplying Eq.~\eqref{eq:14.1.8} on the right with $u_M(\mathbf y)$ or
$v_M(\mathbf y)$, we find then that these wave functions must satisfy the
normalization conditions
\begin{equation}
  \int d^3y\,(u_N^\dagger(\mathbf y)u_M(\mathbf y))
  =\int d^3y\,(v_N^\dagger(\mathbf y)v_M(\mathbf y))=\delta_{NM},
  \tag{14.1.18}\label{eq:14.1.18}
\end{equation}
where $\delta_{NM}$ is a product of Kronecker deltas and momentum space delta
functions, with normalization adapted to that used in defining $\sum_N$, in
such a way that $\sum_N\delta_{NM}=1$. These normalization conditions have
nothing directly to do with any probabilistic interpretation of the Dirac wave
functions, but arise instead from the anticommutation relations
\eqref{eq:14.1.7} for the fields.

Let us now specialize to the case of a pure electrostatic external field with
$\mathbf A=0$. In our standard representation of the Dirac matrices, we have
\[
  \boldsymbol\gamma=i\begin{pmatrix}0&-\boldsymbol\sigma\\
  \boldsymbol\sigma&0\end{pmatrix},\qquad
  i\gamma^0=\beta=\begin{pmatrix}0&\mathbf{1}\\\mathbf{1}&0\end{pmatrix},
\]
where $\boldsymbol\sigma$ is the usual three-vector of $2\times2$ Pauli
matrices, and `$\mathbf{1}$' and `$0$' here are the $2\times2$ unit and zero
matrices. We introduce two-component wave functions $f_N$ and $g_N$ by setting
\begin{equation}
  u_N=\frac1{\sqrt2}\begin{pmatrix}f_N+ig_N\\f_N-ig_N\end{pmatrix}.
  \tag{14.1.19}\label{eq:14.1.19}
\end{equation}
The energy eigenvalue condition \eqref{eq:14.1.10} then takes the form:
\begin{equation}
  (\boldsymbol\sigma\cdot\boldsymbol\nabla)f_N
  =(E_N+e\mathcal A^0+m)g_N,
  \tag{14.1.20}\label{eq:14.1.20}
\end{equation}
\begin{equation}
  (\boldsymbol\sigma\cdot\boldsymbol\nabla)g_N
  =-(E_N+e\mathcal A^0-m)f_N.
  \tag{14.1.21}\label{eq:14.1.21}
\end{equation}
In the non-relativistic case where $e\mathcal A^0r\approx Z\alpha\ll1$ the
binding energy $m-E_N$ is of order $Z^2\alpha^2m$, while the gradient operator
is of the order of $Z\alpha m$, so $g_N$ is smaller than $f_N$ by a factor of
order $Z\alpha$. (To find the positron wave functions $v_N$ we replace $E_N$
everywhere by $-E_N$, so in this case $f_N$ is smaller than $g_N$ by the same
factor.) We shall return to this non-relativistic case at the end of this
section.

Physical states may be classified as even or odd under space inversion:
\begin{equation}
  P\ket{N}=\eta_N\ket{N},
  \tag{14.1.22}\label{eq:14.1.22}
\end{equation}
where $\eta_N$ is a sign factor, $\pm1$. Recall that with the intrinsic parity
of the electron defined to be $+1$, the Dirac field has the space-inversion
property
\[
  P\psi(\mathbf x,t)P^{-1}=\beta\psi(-\mathbf x,t),
\]
so Eqs.~\eqref{eq:14.1.4} and \eqref{eq:14.1.22} show that the Dirac wave
functions satisfy the parity condition
\begin{equation}
  u_N(\mathbf x)=\eta_N\beta u_N(-\mathbf x),
  \tag{14.1.23}\label{eq:14.1.23}
\end{equation}
or equivalently
\begin{equation}
  f_N(\mathbf x)=\eta_Nf_N(-\mathbf x),\qquad
  g_N(\mathbf x)=-\eta_Ng_N(-\mathbf x).
  \tag{14.1.24}\label{eq:14.1.24}
\end{equation}
Note that the parity of the state is the same as the parity of $f_N(\mathbf x)$,
not $g_N(\mathbf x)$.

Where the potential $\mathcal A^0$ is rotationally invariant, the solutions of
the wave equations here may be classified according to their total angular
momentum $j$ and parity $\eta$. For a given $j$, the components $f$ and $g$ may
be expanded in spherical harmonics with orbital angular momentum
$\ell=j+\frac12$ and $\ell=j-\frac12$, but for a definite parity
$\eta=(-1)^{j\mp\frac12}$, Eqs.~\eqref{eq:14.1.24} show that we can only have
$\ell=j\mp\frac12$ in $f$ and $\ell=j\pm\frac12$ in $g$. The usual rules of
angular-momentum addition then show that for a state of total angular momentum
$j$, total angular-momentum $z$-component $\mu$, and parity
$(-1)^{j\mp\frac12}$, the `large' two-component wave function $f$ has the form
\begin{equation}
 f(\mathbf x)=
 \begin{pmatrix}
 C_{j\mp\frac12,\frac12}(j\,\mu;\mu-\frac12,\frac12)
 Y_{j\mp\frac12}^{\mu-\frac12}(\hat{\mathbf x})\\
 C_{j\mp\frac12,\frac12}(j\,\mu;\mu+\frac12,-\frac12)
 Y_{j\mp\frac12}^{\mu+\frac12}(\hat{\mathbf x})
 \end{pmatrix}F(|\mathbf x|),
 \tag{14.1.25}\label{eq:14.1.25}
\end{equation}
where $C$ and $Y$ are the usual Clebsch--Gordan coefficients and spherical
harmonics~\hyperlink{ref:14.2}{[2]}. Also, given any wave function of definite
total angular momentum and parity we can construct another wave function with
the same $j$ and $\mu$ but opposite parity by applying the operator
$\boldsymbol\sigma\cdot\hat{\mathbf x}$, so the `small' components may be put
in the form
\begin{equation}
 g(\mathbf x)=(\boldsymbol\sigma\cdot\hat{\mathbf x})
 \begin{pmatrix}
 C_{j\pm\frac12,\frac12}(j\,\mu;\mu-\frac12,\frac12)
 Y_{j\pm\frac12}^{\mu-\frac12}(\hat{\mathbf x})\\
 C_{j\pm\frac12,\frac12}(j\,\mu;\mu+\frac12,-\frac12)
 Y_{j\pm\frac12}^{\mu+\frac12}(\hat{\mathbf x})
 \end{pmatrix}G(|\mathbf x|).
 \tag{14.1.26}\label{eq:14.1.26}
\end{equation}
It is conventional to define the orbital angular momentum $\ell$ of the state
as the orbital angular momentum of the `large' components $f(\mathbf x)$,
\begin{equation}
  \ell=j\mp\frac12,
  \tag{14.1.27}\label{eq:14.1.27}
\end{equation}
so that the parity is always $(-1)^\ell$.

Inserting Eqs.~\eqref{eq:14.1.25} and \eqref{eq:14.1.26} in
Eqs.~\eqref{eq:14.1.20} and \eqref{eq:14.1.21} yields the coupled differential
equations
\begin{equation}
  \frac{dG}{dr}+\frac{k+1}{r}G+(E+e\mathcal A^0-m)F=0,
  \tag{14.1.28}\label{eq:14.1.28}
\end{equation}
\begin{equation}
  \frac{dF}{dr}-\frac{k-1}{r}F-(E+e\mathcal A^0+m)G=0,
  \tag{14.1.29}\label{eq:14.1.29}
\end{equation}
where for parity $\eta=(-1)^{j\mp\frac12}$,
\begin{equation}
  k=\pm\left(j+\frac12\right).
  \tag{14.1.30}\label{eq:14.1.30}
\end{equation}

Let us now concentrate on the simple Coulomb field \eqref{eq:14.1.2}, for which
$e\mathcal A^0=Z\alpha/r$. The treatment of the Dirac equation in this
case~\hyperlink{ref:14.3}{[3]} is familiar, so it will be summarized briefly
here just for completeness. It is easy to see that the solutions near the
origin go as $r^{s-1}$, with $s^2=k^2-Z^2\alpha^2$. (Note that $k^2\geq1$, so
the exponent $s$ is real for $Z\alpha\leq1$.) We must reject the solutions with
$s<0$ as being inconsistent with the normalization condition
\eqref{eq:14.1.18}. The condition that the wave functions do not blow up for
$r\to\infty$ then fixes the allowed values of the energy eigenvalues:
\begin{equation}
 E_{n,j}=m\left[1+
 \frac{Z^2\alpha^2}{\left[n-j-\frac12+
 \sqrt{(j+\frac12)^2-Z^2\alpha^2}\right]^2}\right]^{-1/2},
 \tag{14.1.31}\label{eq:14.1.31}
\end{equation}
where $n$ is a `principal quantum number' with
\begin{equation}
  j+\frac12\leq n.
  \tag{14.1.32}\label{eq:14.1.32}
\end{equation}
It is noteworthy that these energies do not depend on the parity or $\ell$,
but only on $n$ and $j$. For each $n$ and $j$ there are two solutions,
corresponding to the two signs of $k$ or the two possible parities, except that
for $n=j+\frac12$ we only have $k>0$ and parity $(-1)^{j-\frac12}$, so that
$\ell=j-\frac12$. With Eq.~\eqref{eq:14.1.32}, this is the same as the familiar
non-relativistic restriction that $\ell\leq n-1$.

For light atoms with $Z\alpha\ll1$, Eq.~\eqref{eq:14.1.31} yields the power
series
\begin{equation}
  E=m\left[1-\frac{Z^2\alpha^2}{2n^2}
  +\frac{Z^4\alpha^4}{n^4}\left(\frac38-\frac{n}{2j+1}\right)+\cdots\right].
  \tag{14.1.33}\label{eq:14.1.33}
\end{equation}
The first two terms, of course, just represent the rest energy and the binding
energy as given by the non-relativistic Schr\"odinger equation. The leading
term that depends on $j$ as well as $n$ is the third term, the first
relativistic correction. For $n=1$ there is only one value of the total
angular momentum, $j=\frac12$, and since here $n=j+\frac12$ there is also only
one parity, $(-1)^{j-\frac12}=+1$, corresponding to $\ell=0$. It is therefore
difficult to see the effects of the relativistic corrections in
Eq.~\eqref{eq:14.1.33} in the $n=1$ states of hydrogenic atoms, though as we
shall see in Section 14.3, this has recently become possible. On the other
hand, for $n=2$ we have a $j=\frac12$ state with both parities (i.e.,
$2s_{1/2}$ and $2p_{1/2}$) as well as a $2p_{3/2}$ state with $j=\frac32$ and
negative parity. Eq.~\eqref{eq:14.1.33} gives the splitting between the $p$
states in hydrogen as
\begin{equation}
  E(2p_{3/2})-E(2p_{1/2})=\frac{\alpha^4m_e}{32}
  =4.5283\cdot10^{-5}\ \mathrm{eV}.
  \tag{14.1.34}\label{eq:14.1.34}
\end{equation}
Such relativistic line splitting is known as the \emph{fine structure} of the
atomic state. From the beginning it was known that this prediction is in good
agreement with the observed fine structure. On the other hand, the Dirac
equation does not yield any energy difference in the $2s_{1/2}$ and $2p_{1/2}$
states, so this is a good place to look for the effects of further corrections,
to be considered in Section 14.3.

Before closing this section we shall consider the approximate forms for the
wave functions and matrix elements in the non-relativistic case for a general
electrostatic potential $\mathcal A^0$. (For a Coulomb potential, this is the
limit $Z\alpha\ll1$.) Since here $E_N+m\simeq2m\gg|e\mathcal A^0|$, the
`small' components of the electron wave function are given approximately in
terms of the large components by
\begin{equation}
  g_N\simeq(\boldsymbol\sigma\cdot\boldsymbol\nabla)f_N/2m_e.
  \tag{14.1.35}\label{eq:14.1.35}
\end{equation}
Eq.~\eqref{eq:14.1.21} then becomes just the non-relativistic Schr\"odinger
equation
\begin{equation}
  \left[-\frac{\boldsymbol\nabla^2}{2m_e}-e\mathcal A^0\right]f_N
  \simeq(E_N-m)f_N.
  \tag{14.1.36}\label{eq:14.1.36}
\end{equation}
Since there is no longer any coupling between spin and orbital degrees of
freedom in the equation for $f_N$, we may find a complete set of solutions of
this equation in the form
\[
  f_N=\chi_N\psi_N(\mathbf x),
\]
where $\chi_N$ is a two-component constant spinor, and $\psi_N(\mathbf x)$ is
an ordinary one-component solution of the Schr\"odinger equation. However, we
often work with states that have definite values of the total angular momentum
$j$, for which $f_N$ is (for non-vanishing orbital angular momentum) a sum of
such terms.

In the non-relativistic approximation, the four-component Dirac wave function
takes the form
\begin{equation}
  u_N\simeq\frac1{\sqrt2}
  \begin{bmatrix}(1+i\boldsymbol\sigma\cdot\boldsymbol\nabla/2m)f_N\\
  (1-i\boldsymbol\sigma\cdot\boldsymbol\nabla/2m)f_N\end{bmatrix}
  \tag{14.1.37}\label{eq:14.1.37}
\end{equation}
and Eq.~\eqref{eq:14.1.18} gives the normalization condition
\begin{equation}
  \int d^3x\,(f_N^\dagger f_M)\simeq\delta_{NM}-\frac14(v^2)_{NM},
  \tag{14.1.38}\label{eq:14.1.38}
\end{equation}
where
\[
  (v^2)_{NM}\equiv-\frac1{m_e^2}\int d^3x\,
  f_N^\dagger(\mathbf x)\boldsymbol\nabla^2f_M(\mathbf x).
\]
In relating matrix elements in an external field to free-particle matrix
elements, it is useful to note that the momentum space wave function in an
energy eigenstate $N$ may be written
\begin{equation}
  u_N(\mathbf p)\equiv(2\pi)^{-3/2}\int d^3p\,
  e^{-i\mathbf p\cdot\mathbf x}u_N(\mathbf x)
  \simeq\sum_\sigma u(\mathbf p,\sigma)[f_N(\mathbf p)]_\sigma,
  \tag{14.1.39}\label{eq:14.1.39}
\end{equation}
where $u(\mathbf p,\sigma)$ is the free-particle Dirac spinor
\[
 u(\mathbf p,\sigma)\simeq\frac1{\sqrt2}
 \begin{bmatrix}(1-\mathbf p\cdot\boldsymbol\sigma/2m_e)\chi_\sigma\\
 (1+\mathbf p\cdot\boldsymbol\sigma/2m_e)\chi_\sigma\end{bmatrix},
 \qquad
 \chi_{+\frac12}\equiv\begin{pmatrix}1\\0\end{pmatrix},\quad
 \chi_{-\frac12}\equiv\begin{pmatrix}0\\1\end{pmatrix},
\]
and $f_N(\mathbf p)$ is the Fourier transform of the two-component
Schr\"odinger wave function
\[
  f_N(\mathbf p)\equiv(2\pi)^{-3/2}\int d^3p\,
  e^{-i\mathbf p\cdot\mathbf x}f_N(\mathbf x).
\]

In closing, as an aid in calculating the effects of various perturbations, we
note that the leading terms in the electron matrix elements of the sixteen
independent $4\times4$ matrices are
\begin{equation}
 (\bar u_Mu_N)\simeq(f_M^\dagger f_N)
 -\frac1{4m_e^2}(\boldsymbol\nabla f_M^\dagger\cdot\boldsymbol\sigma\,
 \boldsymbol\sigma\cdot\boldsymbol\nabla f_N),
 \tag{14.1.40}\label{eq:14.1.40}
\end{equation}
\begin{equation}
 i(\bar u_M\gamma^0u_N)\simeq(f_M^\dagger f_N)
 +\frac1{4m_e^2}(\boldsymbol\nabla f_M^\dagger\cdot\boldsymbol\sigma\,
 \boldsymbol\sigma\cdot\boldsymbol\nabla f_N),
 \tag{14.1.41}\label{eq:14.1.41}
\end{equation}
\begin{equation}
 (\bar u_M\gamma^iu_N)\simeq\frac1{2m_e}
 \left[(\boldsymbol\nabla f_M^\dagger\cdot\boldsymbol\sigma\,\sigma_i f_N)
 -(f_M^\dagger\sigma_i\boldsymbol\sigma\cdot\boldsymbol\nabla f_N)\right],
 \tag{14.1.42}\label{eq:14.1.42}
\end{equation}
\begin{equation}
 (\bar u_M[\gamma^0,\gamma^i]u_N)\simeq\frac{i}{m_e}
 \left[(\boldsymbol\nabla f_M^\dagger\cdot\boldsymbol\sigma\,\sigma_i f_N)
 +(f_M^\dagger\sigma_i\boldsymbol\sigma\cdot\boldsymbol\nabla f_N)\right],
 \tag{14.1.43}\label{eq:14.1.43}
\end{equation}
\begin{equation}
 (\bar u_M[\gamma^i,\gamma^j]u_N)\simeq
 2i\epsilon_{ijk}(f_M^\dagger\sigma_kf_N),
 \tag{14.1.44}\label{eq:14.1.44}
\end{equation}
\begin{equation}
 (\bar u_M\gamma_5\gamma^iu_N)\simeq-i(f_M^\dagger\sigma_if_N),
 \tag{14.1.45}\label{eq:14.1.45}
\end{equation}
\begin{equation}
 (\bar u_M\gamma_5\gamma^0u_N)\simeq\frac1{2m_e}
 \left[(\boldsymbol\nabla f_M^\dagger\cdot\boldsymbol\sigma f_N)
 -(f_M^\dagger\boldsymbol\sigma\cdot\boldsymbol\nabla f_N)\right],
 \tag{14.1.46}\label{eq:14.1.46}
\end{equation}
\begin{equation}
 (\bar u_M\gamma_5u_N)\simeq\frac{i}{2m_e}
 \left[(\boldsymbol\nabla f_M^\dagger\cdot\boldsymbol\sigma f_N)
 +(f_M^\dagger\boldsymbol\sigma\cdot\boldsymbol\nabla f_N)\right].
 \tag{14.1.47}\label{eq:14.1.47}
\end{equation}
